% =============================================================================
%  Ferran Alía Castillejos: curriculum vitae
%
%  Content only. Every design decision lives in cvstyle.sty, which carries the
%  same tokens as the website's src/styles/global.css.
%
%  Build:  ./build.sh        (xelatex, twice, output copied to public/)
%
%  Facts here must agree with src/data/profile.js, src/data/writing.js and
%  src/content/projects/. If you change one, change the other.
% =============================================================================

\documentclass[10pt,a4paper]{article}
\usepackage{cvstyle}

\setcvfooter{Ferran Alía Castillejos}

\begin{document}

\cvheader
  {Ferran Alía Castillejos}
  {Data science and physical engineering, UPC Barcelona}
  {%
    \href{mailto:ferran.alia@estudiantat.upc.edu}{ferran.alia@estudiantat.upc.edu}%
    \cvdot
    \href{https://feal-ca.github.io}{feal-ca.github.io}%
    \cvdot
    \href{https://github.com/Feal-ca}{github.com/Feal-ca}%
    \cvdot
    \href{https://www.linkedin.com/in/ferran-alca/}{linkedin.com/in/ferran-alca}%
  }

% -----------------------------------------------------------------------------
\begin{cvsection}{Now}
I simulate microscope optics with the Physics of Life group at TU Dresden
for the summer, then go back to Barcelona for my fourth year at UPC. I want
to work where physical modeling and machine learning meet.
\end{cvsection}

% -----------------------------------------------------------------------------
\begin{cvsection}{Experience}

\cvitem{Since 2026}{Contributor, \href{https://towardsdatascience.com/author/ferran-alia/}{Towards Data Science}}
  {Long-form articles since March 2026, on simulation, high-performance
   computing and machine learning. Most of them come out of the projects
   below.}

\cvitem{2026}{Research intern, Physics of Life, TU Dresden}
  {Simulating electromagnetic propagation through specimen and objective, as
   virtual instrumentation for learned brightfield imaging.}

\cvitem{Jun\,--\,Aug 2025}{Machine learning intern, MLCode}
  {Built an automated framework for evaluating and benchmarking
   retrieval-augmented generation systems.}

\cvitem{Since 2022}{Volunteer mentor}
  {Computer science and robotics workshops: Linux, 3D printing, programming.}

\cvitemlast{2022\,--\,2023}{Robotics teacher, Punt Multimèdia}
  {Taught Arduino, micro:bit, and 3D printing and modeling.}

\end{cvsection}

% -----------------------------------------------------------------------------
\begin{cvsection}{Education}

\cvitem{Since 2023}{BSc Data Science \& Engineering, UPC}
  {Entering the fourth year of five. 8.75/10 average across both degrees.}

\cvitemlast{Since 2023}{BSc Physical Engineering, UPC}
  {Taken in parallel with the data science degree rather than after it, on
   the CFIS dual-degree program.}

\end{cvsection}

% -----------------------------------------------------------------------------
\begin{cvsection}{Selected projects}

\cvproject{2026}{f1-frontwing}{Aerodynamic optimization of a Formula 1 front wing}
  {A surrogate model searched front-wing geometries on a fixed budget of
   solver runs rather than calling CFD at every step. The wing came out about
   15\% better on lift-to-drag.%
   \cvstack{Python \textperiodcentered\ OpenFOAM \textperiodcentered\ SLURM \textperiodcentered\ MareNostrum V}}

\cvproject{2026}{vortex-pinn}{A physics-informed network for vortex shedding}
  {A PINN with the Navier-Stokes residual inside its loss, trained against a
   high-fidelity OpenFOAM run, reproducing Kármán vortex shedding. A model
   fit to the data alone settles into a steady wake instead.%
   \cvstack{Python \textperiodcentered\ PyTorch \textperiodcentered\ OpenFOAM}}

\cvproject{2026}{llm-persona}{Persona adaptation in language models}
  {A system prompt against three formats of fine-tuning data on
   Qwen3-4B-Instruct, all with LoRA at r=16, so the differences belong to the
   data.%
   \cvstack{Python \textperiodcentered\ PyTorch \textperiodcentered\ Transformers \textperiodcentered\ LoRA}}

\cvproject{2025}{assisted-birth}{Predicting fetal distress from cardiotocography}
  {Built with Hospital Sant Joan de Déu in a team of nine. Separate encoders
   over CTG traces, tabular clinical context and ultrasound, at around 0.78
   AUC.%
   \cvstack{Python \textperiodcentered\ PyTorch}}

\cvproject{2025}{n-body}{Parallel N-body simulation}
  {Direct O(n²) summation and a Barnes-Hut O(n log n) tree code in C++
   with OpenMP. At N = 10,000 direct summation ran about 33× faster on
   112 threads than on one; the tree code never got past roughly 7×.%
   \cvstack{C++ \textperiodcentered\ OpenMP}}

\cvproject{2025}{lattice-boltzmann}{Parallel lattice Boltzmann solver}
  {A D2Q9 solver with BGK collision and a Zou-He inlet in around 200 lines of
   C++: flow past a cylinder on a 400×400 lattice. Memory-bound, so
   how you move the data dominates the scaling.%
   \cvstack{C++ \textperiodcentered\ OpenMP \textperiodcentered\ MareNostrum 5}}

\cvproject{2025}{bird-flight}{Fixed-wing bird flight simulation}
  {A 3D Navier-Stokes solver written from scratch in NumPy, run on a bird's
   wing modeled in Blender and checked against OpenFOAM's icoFoam.%
   \cvstack{Python \textperiodcentered\ NumPy \textperiodcentered\ OpenFOAM \textperiodcentered\ Blender}}

\cvprojectlast{2024}{fermi-dirac}{Fermi-Dirac statistics by Monte Carlo}
  {A Markov chain over the microstates of a fermion system, recovering the
   Fermi-Dirac occupancy numerically rather than analytically.%
   \cvstack{Python \textperiodcentered\ NumPy}}

\cvnote{Figures are schematic illustrations of each problem, not results.}

\end{cvsection}

% -----------------------------------------------------------------------------
\begin{cvsection}{Skills}

\cvchiprow{Languages}{%
  \cvchip{Python}\cvchip{C++}\cvchip{R}\cvchip{Haskell}\cvchip{MATLAB}}

\cvchiprow{Scientific}{%
  \cvchip{OpenFOAM}\cvchip{OpenMP}\cvchip{NumPy}\cvchip{PyTorch}%
  \cvchip{Blender}\cvchip{CAD}}

\cvchiprow{Methods}{%
  \cvchip{CFD}\cvchip{HPC \& SLURM}\cvchip{Physics-informed ML}%
  \cvchip{Monte Carlo}\cvchip{Surrogate modeling}}

\cvchiprow{Spoken}{%
  \cvchip{Catalan (native)}\cvchip{Spanish (native)}%
  \cvchip{English (professional)}\cvchip{Arabic (beginner)}}

\end{cvsection}

% -----------------------------------------------------------------------------
\begin{cvsection}{Awards}

\cvlistrow{2024, 2025}{Participant, Datathon FME}
\cvlistrow{2025}{Physics summer school, University of Ljubljana}
\cvlistrow{Since 2023}{Cum laude in several subjects}
\cvlistrow{2023}{Admitted to CFIS on a grant, to read the two degrees at
  once. It takes a small cohort each year by its own entrance exam in
  mathematics and physics}
\cvlistrow{2019\,--\,2023}{Graduated from high school with honors, and
  represented it at the Barcelona \emph{Mostra de Recerca Jove}, after UPC's
  summer courses in algorithms and programming}

\end{cvsection}

\cvtail
\label{cv:last}
\end{document}
